\section{Question 3 --- Products, Ores and Sensitivity}

\begin{questionbox}
A company produces three products \(A\), \(B\) and \(C\). These products require three ores
\(O_1\), \(O_2\) and \(O_3\). The maximum quantities of the ores \(O_1\), \(O_2\) and \(O_3\)
available are \(22\) tonnes, \(14\) tonnes and \(14\) tonnes, respectively. For one tonne of
each of these products, the ore requirements are:

\begin{center}
\begin{tabular}{c|ccc}
 & \(A\) & \(B\) & \(C\) \\
\hline
\(O_1\) & 3 & --- & 3 \\
\(O_2\) & 1 & 2 & 3 \\
\(O_3\) & 3 & 2 & 3 \\
\hline
Profit per tonne & 1000 & 4000 & 5000 \\
\end{tabular}
\end{center}

\begin{enumerate}
  \item[(a)] How many tonnes of each product should the company produce in order to
        maximize its profits?
  \item[(b)] Determine the range of feasibility for the right-hand side value of ore \(O_2\).
  \item[(c)] Determine the range of optimality for the profit coefficient of product \(A\).
\end{enumerate}
\end{questionbox}

\begin{notebox}
The dash for \(B\) under \(O_1\) means product \(B\) uses \textbf{no} ore \(O_1\)
(coefficient \(0\)).
\end{notebox}

% ----------------------------------------------------------------------
\PartHeading{(a) Optimal production plan}
% ----------------------------------------------------------------------

\begin{ideabox}
Formulate the LP, put it in standard form, then run the simplex algorithm until all
\(C_j-Z_j\le 0\).
\end{ideabox}

\Step{Decision variables}
Let \(x_A,x_B,x_C\) be tonnes of products \(A\), \(B\), \(C\).

\Step{LP model}
\begin{align*}
  \max\; Z &= 1000x_A + 4000x_B + 5000x_C \\
  \text{s.t.}\quad
  3x_A + 3x_C &\le 22 && (O_1),\\
  x_A + 2x_B + 3x_C &\le 14 && (O_2),\\
  3x_A + 2x_B + 3x_C &\le 14 && (O_3),\\
  x_A,x_B,x_C &\ge 0.
\end{align*}

\Step{Standard form}
\begin{align*}
  3x_A + 0\cdot x_B + 3x_C + s_1 &= 22,\\
  x_A + 2x_B + 3x_C + s_2 &= 14,\\
  3x_A + 2x_B + 3x_C + s_3 &= 14,\\
  Z - 1000x_A - 4000x_B - 5000x_C &= 0.
\end{align*}

\Step{Initial tableau}
\begin{center}
\renewcommand{\arraystretch}{1.15}
\begin{tabular}{c|c|ccc|ccc|c}
\toprule
\(C_B\) & Basis & \(x_A\) & \(x_B\) & \(x_C\) & \(s_1\) & \(s_2\) & \(s_3\) & RHS \\
\midrule
0 & \(s_1\) & 3 & 0 & 3 & 1 & 0 & 0 & 22 \\
0 & \(s_2\) & 1 & 2 & 3 & 0 & 1 & 0 & 14 \\
0 & \(s_3\) & 3 & 2 & 3 & 0 & 0 & 1 & 14 \\
\midrule
 & \(C_j-Z_j\) & 1000 & 4000 & \textbf{5000} & 0 & 0 & 0 & \(Z=0\) \\
\bottomrule
\end{tabular}
\end{center}

Enter \(x_C\) (largest indicator \(5000\)). Ratios:
\[
  \tfrac{22}{3},\qquad \tfrac{14}{3},\qquad \tfrac{14}{3}.
\]
There is a tie between \(s_2\) and \(s_3\); we let \(s_2\) leave (either choice is valid).
Pivot \(=3\).

\Step{Iteration 1}
After pivoting \(x_C\) into the basis:
\begin{center}
\renewcommand{\arraystretch}{1.15}
\begin{tabular}{c|c|ccc|ccc|c}
\toprule
\(C_B\) & Basis & \(x_A\) & \(x_B\) & \(x_C\) & \(s_1\) & \(s_2\) & \(s_3\) & RHS \\
\midrule
0 & \(s_1\) & 2 & \(-2\) & 0 & 1 & \(-1\) & 0 & 8 \\
5000 & \(x_C\) & \(\tfrac13\) & \(\tfrac23\) & 1 & 0 & \(\tfrac13\) & 0 & \(\tfrac{14}{3}\) \\
0 & \(s_3\) & 2 & 0 & 0 & 0 & \(-1\) & 1 & 0 \\
\midrule
 & \(C_j-Z_j\) & \(-\tfrac{2000}{3}\) & \(\mathbf{\tfrac{2000}{3}}\) & 0 & 0 & \(-\tfrac{5000}{3}\) & 0 & \(Z=\tfrac{70000}{3}\) \\
\bottomrule
\end{tabular}
\end{center}

How the indicators were obtained:
\begin{align*}
  C_A-Z_A &= 1000 - 5000\cdot\tfrac13 = -\tfrac{2000}{3},\\
  C_B-Z_B &= 4000 - 5000\cdot\tfrac23 = \tfrac{2000}{3}>0.
\end{align*}
So enter \(x_B\). In the \(x_B\) column the only useful positive pivot candidate is
\(\frac{2}{3}\) in the \(x_C\)-row (the \(s_1\) entry is \(-2<0\); the \(s_3\) entry is \(0\)).
Ratio \(=\frac{14/3}{2/3}=7\). Thus \(x_C\) leaves.

\Step{Iteration 2 --- optimal tableau}
\begin{center}
\renewcommand{\arraystretch}{1.2}
\begin{tabular}{c|c|ccc|ccc|c}
\toprule
\(C_B\) & Basis & \(x_A\) & \(x_B\) & \(x_C\) & \(s_1\) & \(s_2\) & \(s_3\) & RHS \\
\midrule
0 & \(s_1\) & 3 & 0 & 3 & 1 & 0 & 0 & 22 \\
4000 & \(x_B\) & \(\tfrac12\) & 1 & \(\tfrac32\) & 0 & \(\tfrac12\) & 0 & 7 \\
0 & \(s_3\) & 2 & 0 & 0 & 0 & \(-1\) & 1 & 0 \\
\midrule
 & \(C_j-Z_j\) & \(-1000\) & 0 & \(-1000\) & 0 & \(-2000\) & 0 & \(Z=28000\) \\
\bottomrule
\end{tabular}
\end{center}

(Equivalently, the optimality row may be written with signs flipped as \(Z_j-C_j\); what
matters is that no non-basic variable improves \(Z\).)

All improvement indicators for non-basics are non-positive \(\Rightarrow\) optimal.

\begin{resultbox}[title={Part (a)}]
Produce
\[
  x_A^*=0,\qquad x_B^*=7,\qquad x_C^*=0
\]
tonnes, with maximum profit
\[
  Z^*=28\,000.
\]
Ore \(O_1\) is unused (\(s_1^*=22\)); ores \(O_2\) and \(O_3\) are fully used.
\end{resultbox}

\begin{warningbox}
The optimal tableau is \emph{degenerate} (\(s_3^*=0\) while basic). That is normal here:
when \(x_A=x_C=0\), constraints \(O_2\) and \(O_3\) both reduce to \(2x_B\le 14\), so they
become the same restriction.
\end{warningbox}

% ----------------------------------------------------------------------
\PartHeading{(b) Range of feasibility for the RHS of \(O_2\)}
% ----------------------------------------------------------------------

\begin{ideabox}
Replace the \(O_2\) right-hand side \(14\) by \(14+\Delta\). The current basis stays feasible
as long as every basic RHS stays non-negative. Those new RHS values equal
\[
  \bar{b} + \Delta\cdot\bigl(B^{-1}e_2\bigr),
\]
and the column \(B^{-1}e_2\) is exactly the final-tableau column of the \(O_2\) slack \(s_2\).
\end{ideabox}

\Step{Read \(B^{-1}e_2\) from the optimal tableau}
Basic variables (in row order): \(s_1\), \(x_B\), \(s_3\). The \(s_2\) column is
\[
  B^{-1}e_2 =
  \begin{pmatrix} 0 \\ \tfrac12 \\ -1 \end{pmatrix}.
\]
Current basic RHS: \(\bar{b}=(22,\;7,\;0)^\top\).

\Step{Impose non-negativity after the change}
\[
\begin{pmatrix} 22 \\ 7 \\ 0 \end{pmatrix}
+
\Delta
\begin{pmatrix} 0 \\ \tfrac12 \\ -1 \end{pmatrix}
\ge
\begin{pmatrix} 0 \\ 0 \\ 0 \end{pmatrix}
\quad\Longleftrightarrow\quad
\begin{cases}
  22 \ge 0 & \text{(always)},\\
  7+\tfrac12\Delta \ge 0 & \Rightarrow\ \Delta\ge -14,\\
  0-\Delta \ge 0 & \Rightarrow\ \Delta\le 0.
\end{cases}
\]
Hence
\[
  -14\le\Delta\le 0
  \qquad\text{i.e.}\qquad
  0\le b_2\le 14.
\]

\begin{resultbox}[title={Part (b)}]
Range of feasibility for the RHS of ore \(O_2\):
\[
  \boxed{\,0\le b_{O_2}\le 14\,}.
\]
(Increasing \(O_2\) above \(14\) cannot help while \(O_3\) is still capped at \(14\), because
\(B\) also needs \(O_3\). Decreasing \(O_2\) below \(0\) is meaningless; at \(0\) we get
\(x_B=0\).)
\end{resultbox}

% ----------------------------------------------------------------------
\PartHeading{(c) Range of optimality for the profit of \(A\)}
% ----------------------------------------------------------------------

\begin{ideabox}
Product \(A\) is \emph{non-basic} at the optimum. Changing its profit coefficient only
changes \emph{its own} reduced cost. The basis (and therefore the production plan)
stays optimal as long as \(A\)'s reduced cost does not become attractive enough to enter.
\end{ideabox}

\Step{Current reduced cost of \(A\)}
From the optimal tableau,
\[
  \bar{c}_A = C_A - Z_A = -1000
\]
(or equivalently \(Z_A-C_A=+1000\)). So bringing one tonne of \(A\) into the plan would
\emph{worsen} profit by \(1000\).

\Step{Allowable change \(\delta\) in \(C_A\)}
Write the new profit as \(C_A'=1000+\delta\). Then the new reduced cost is
\[
  \bar{c}_A' = (1000+\delta)-Z_A = \bar{c}_A+\delta = -1000+\delta.
\]
For maximisation we need \(\bar{c}_A'\le 0\):
\[
  -1000+\delta\le 0 \quad\Rightarrow\quad \delta\le 1000.
\]
Decreasing \(C_A\) (negative \(\delta\)) only makes \(A\) less attractive, so there is no
finite lower bound from optimality.

Therefore
\[
  C_A' = 1000+\delta \le 2000,
  \qquad\text{i.e.}\qquad
  C_A\in(-\infty,\,2000].
\]

\begin{resultbox}[title={Part (c)}]
Range of optimality for the profit coefficient of product \(A\):
\[
  \boxed{\,C_A\le 2000\,}
\]
(with the usual implicit \(C_A\) free to fall arbitrarily). As long as \(A\)'s unit profit
does not exceed \(2000\), the optimal basis stays \(\{s_1,x_B,s_3\}\) and the company
still produces \(7\) tonnes of \(B\) only.
\end{resultbox}
